\chapter{Boolean Analysis — Restriction Card Tail}
%%% generated by the blueprint pipeline from TCSlib.BooleanAnalysis.LMN.RestrictionCardTail
\section{Overview}
Under a Bernoulli($p$)-random restriction with free-coordinate set $J$, the count
$\abs{U \cap J}$ of free coordinates inside a fixed set $U$ behaves like a
$\mathrm{Binomial}(\abs{U}, p)$ variable. This module computes its first and second moments
from the marginal $\Pr[T \subseteq J] = p^{\abs{T}}$ and derives the elementary Chebyshev
lower-tail bounds $\Pr[\abs{U \cap J} < k] \le 3/(4k)$ and $\Pr[\abs{U \cap J} \ge k] \ge 1/4$
used by the restriction $\Rightarrow$ Fourier-concentration transfer (O'Donnell Lemma 4.21),
together with the complement rule for Bernoulli-restriction probabilities.

\section{Declarations}

\begin{lemma}[Indicator of a subset lying among free variables as a coordinate product]
\lean{RestrictionFourier.indicator_subset_eq_prod}
\label{RestrictionFourier.indicator_subset_eq_prod}
\leanok
\uses{RestrictionFourier.mem_freeVars, SwitchingLemma2.Restriction, SwitchingLemma2.Restriction.freeVars}
\difficulty{4}
Let $\rho$ be a restriction on $n$ variables — a map assigning each index $i \in
\{0,\dots,n-1\}$ either a fixed Boolean value or the free mark — and let $T \subseteq
\{0,\dots,n-1\}$. Then the indicator of the event that every index in $T$ is free under
$\rho$ factors as a product over all $n$ coordinates: \[ \mathbf{1}\bigl[\,T \subseteq
\{i : \rho(i)\ \text{is free}\}\,\bigr] \;=\; \prod_{i=0}^{n-1} c_i, \] where the factor
$c_i$ equals $\mathbf{1}[\rho(i)\ \text{is free}]$ when $i \in T$ and equals $1$
otherwise.
\end{lemma}

\begin{theorem}[Free-set marginal of a Bernoulli restriction]
\lean{RestrictionFourier.bernoulliRestrProb_subset_freeVars}
\label{RestrictionFourier.bernoulliRestrProb_subset_freeVars}
\leanok
\uses{LMN.varWeight, RestrictionFourier.indicator_subset_eq_prod, RestrictionFourier.sum_bernoulli_prod, SwitchingLemma2.Restriction.freeVars, SwitchingLemma2.bernoulliRestrProb}
\difficulty{4}
Fix $n \in \bbn$, a real parameter $p$, and a subset $T \subseteq \{1, \dots, n\}$ of
the coordinates. Draw a Bernoulli($p$) random restriction $\rho$ on $n$ variables, and
write $J$ for its set of free variables (the coordinates $\rho$ leaves unset). Then the
probability that every coordinate of $T$ is left free is
\[
  \Pr\bigl[\,T \subseteq J\,\bigr] \;=\; p^{\abs{T}},
\]
that is, the total Bernoulli weight of those restrictions whose free variables include
$T$ equals $p^{\abs{T}}$.
\end{theorem}

\begin{lemma}[Free-variable count as a sum of indicators]
\lean{RestrictionFourier.card_inter_eq_sum}
\label{RestrictionFourier.card_inter_eq_sum}
\leanok
\uses{RestrictionFourier.mem_freeVars, SwitchingLemma2.Restriction, SwitchingLemma2.Restriction.freeVars}
\difficulty{3}
Let $\rho$ be a restriction on $n$ variables, and let $U \subseteq \{1,\dots,n\}$ be a
set of indices. Then the number of indices in $U$ at which $\rho$ is free (unset),
viewed as a real number, equals
\[
\abs{U \cap J} = \sum_{i \in U} \mathbf{1}[\rho(i)\ \text{is free}],
\]
where $J$ is the set of free variables of $\rho$ and each summand is $1$ when $\rho(i)$
is unset and $0$ otherwise.
\end{lemma}

\begin{lemma}[Total Bernoulli weight of restrictions leaving a coordinate free]
\lean{RestrictionFourier.expectation_free_indicator}
\label{RestrictionFourier.expectation_free_indicator}
\leanok
\uses{RestrictionFourier.bernoulliRestrProb_subset_freeVars, RestrictionFourier.mem_freeVars, SwitchingLemma2.Restriction, SwitchingLemma2.Restriction.freeVars, SwitchingLemma2.bernoulliRestrProb, SwitchingLemma2.bernoulliRestrWeight}
\difficulty{4}
Fix a real parameter $p$ and a coordinate $i$ among the $n$ variables. Summing the
Bernoulli weight of $\rho$ over all restrictions $\rho$ on $n$ variables that leave $i$
free — that is, those for which $i$ is a free variable of $\rho$ — gives
\[
\sum_{\rho \,:\, i \text{ free in } \rho}
p^{\,|F(\rho)|}\left(\tfrac{1-p}{2}\right)^{\,n-|F(\rho)|} = p,
\]
where $F(\rho)$ denotes the set of free variables of $\rho$.
\end{lemma}

\begin{lemma}[Probability that two fixed coordinates are both free]
\lean{RestrictionFourier.expectation_free_indicator_pair}
\label{RestrictionFourier.expectation_free_indicator_pair}
\leanok
\uses{RestrictionFourier.bernoulliRestrProb_subset_freeVars, RestrictionFourier.mem_freeVars, SwitchingLemma2.Restriction, SwitchingLemma2.Restriction.freeVars, SwitchingLemma2.bernoulliRestrProb, SwitchingLemma2.bernoulliRestrWeight}
\difficulty{4}
Fix $n$ and two distinct variables $i \neq j$ among the $n$ variables, and let $p \in
\bbr$. Write $w_p(\rho)$ for the Bernoulli weight of a restriction $\rho$ on $n$
variables, namely $p^{k}\left(\tfrac{1-p}{2}\right)^{n-k}$ where $k$ is the number of
free variables of $\rho$. Then, summing over all such restrictions, the total weight
carried by those $\rho$ that leave both $i$ and $j$ free equals $p^2$; that is,
\[
\sum_{\rho} w_p(\rho)\,\mathbf{1}[\rho(i)\text{ free}]\,\mathbf{1}[\rho(j)\text{ free}]
= p^2.
\]
\end{lemma}

\begin{lemma}[First moment of the free-variable count under a Bernoulli restriction]
\lean{RestrictionFourier.expectation_card_inter}
\label{RestrictionFourier.expectation_card_inter}
\leanok
\uses{RestrictionFourier.card_inter_eq_sum, RestrictionFourier.expectation_free_indicator, SwitchingLemma2.Restriction, SwitchingLemma2.Restriction.freeVars, SwitchingLemma2.bernoulliRestrWeight}
\difficulty{3}
Fix $n$ and a real number $p$, and let $\rho$ range over all restrictions on the $n$
variables, each weighted by its Bernoulli weight $w_p(\rho)$. Writing $J(\rho)$ for the
set of free variables of $\rho$, then for every set $U$ of variables,
\[
  \sum_{\rho} w_p(\rho)\,\abs{U \cap J(\rho)} \;=\; p\,\abs{U}.
\]
Equivalently, when $0 \le p \le 1$ so that $w_p$ is a probability distribution, the
expected number of free variables landing in $U$ is $p\,\abs{U}$.
\end{lemma}

\begin{lemma}[Second moment of the free-coordinate count]
\lean{RestrictionFourier.expectation_card_inter_sq}
\label{RestrictionFourier.expectation_card_inter_sq}
\leanok
\uses{RestrictionFourier.card_inter_eq_sum, RestrictionFourier.expectation_free_indicator, RestrictionFourier.expectation_free_indicator_pair, SwitchingLemma2.Restriction, SwitchingLemma2.Restriction.freeVars, SwitchingLemma2.bernoulliRestrWeight}
\difficulty{5}
Fix $n \in \bbn$, a parameter $p \in \bbr$, and a set of coordinates $U \subseteq
\mathrm{Fin}\,n$. For a restriction $\rho$ on $n$ variables write $J(\rho)$ for its set
of free variables. Summing over all such restrictions, each weighted by its Bernoulli
weight, one has
\[
\sum_{\rho} p^{\abs{J(\rho)}}\left(\tfrac{1-p}{2}\right)^{\,n - \abs{J(\rho)}} \abs{U
\cap J(\rho)}^2
  \;=\; p\,\abs{U} + p^2\bigl(\abs{U}^2 - \abs{U}\bigr).
\]
\end{lemma}

\begin{lemma}[Variance of the free-coordinate count]
\lean{RestrictionFourier.variance_card_inter}
\label{RestrictionFourier.variance_card_inter}
\leanok
\uses{RestrictionFourier.expectation_card_inter, RestrictionFourier.expectation_card_inter_sq, SwitchingLemma2.Restriction, SwitchingLemma2.Restriction.freeVars, SwitchingLemma2.bernoulliRestrWeight, SwitchingLemma2.bernoulliRestrWeight_sum_one}
\difficulty{4}
Let $0 \le p \le 1$, and let $U \subseteq \mathrm{Fin}\,n$ be a set of coordinates. For
a restriction $\rho$ on $n$ variables, write $w_p(\rho)$ for its Bernoulli weight and
$F(\rho)$ for its set of free variables. Then the $w_p$-weighted average, over all
restrictions $\rho$ on $n$ variables, of the squared deviation of $\abs{U \cap F(\rho)}$
from $p\abs{U}$ satisfies
\[
\sum_{\rho}\, w_p(\rho)\,\bigl(p\abs{U} - \abs{U \cap F(\rho)}\bigr)^2 \;=\;
p\,\abs{U}\,(1 - p).
\]
\end{lemma}

\begin{theorem}[Lower-tail bound for free coordinates in a set]
\lean{RestrictionFourier.bernoulliRestrProb_card_inter_lt}
\label{RestrictionFourier.bernoulliRestrProb_card_inter_lt}
\leanok
\uses{RestrictionFourier.variance_card_inter, SwitchingLemma2.Restriction, SwitchingLemma2.Restriction.freeVars, SwitchingLemma2.bernoulliRestrProb, SwitchingLemma2.bernoulliRestrWeight, SwitchingLemma2.bernoulliRestrWeight_nonneg'}
\difficulty{6}
Let $p \in \bbr$ satisfy $0 \le p \le 1$, let $U \subseteq \mathrm{Fin}\,n$ be a finite
set of coordinates, and let $k$ be a positive integer with $3k \le p\,\abs{U}$. Writing
$J$ for the set of free variables of a Bernoulli($p$)-random restriction $\rho$ on $n$
variables, the probability that fewer than $k$ coordinates of $U$ remain free satisfies
\[
  \Pr\bigl[\,\abs{U \cap J} < k\,\bigr] \;\le\; \frac{3}{4k}.
\]
\end{theorem}

\begin{lemma}[Complement rule for the Bernoulli restriction measure]
\lean{RestrictionFourier.bernoulliRestrProb_not}
\label{RestrictionFourier.bernoulliRestrProb_not}
\leanok
\uses{SwitchingLemma2.Restriction, SwitchingLemma2.bernoulliRestrProb, SwitchingLemma2.bernoulliRestrWeight_sum_one}
\difficulty{3}
Fix $n$ and a parameter $p$ with $0 \le p \le 1$, and consider the Bernoulli($p$) random
restriction on $n$ variables, under which each restriction $\rho$ receives weight
$p^{k}\,\bigl(\tfrac{1-p}{2}\bigr)^{\,n-k}$, where $k$ is the number of free variables
of $\rho$. For any decidable predicate $E$ on restrictions, writing $\Pr[E]$ for the
total weight of the restrictions satisfying $E$, the probability that $E$ fails equals
one minus the probability that $E$ holds:
\[
  \Pr[\lnot E] = 1 - \Pr[E].
\]
\end{lemma}

\begin{theorem}[Lower bound on the free coordinates within a set]
\lean{RestrictionFourier.bernoulliRestrProb_card_inter_ge}
\label{RestrictionFourier.bernoulliRestrProb_card_inter_ge}
\leanok
\uses{RestrictionFourier.bernoulliRestrProb_card_inter_lt, RestrictionFourier.bernoulliRestrProb_not, SwitchingLemma2.Restriction.freeVars, SwitchingLemma2.bernoulliRestrProb}
\difficulty{4}
Fix $n$ and a parameter $p$ with $0 \le p \le 1$, let $U \subseteq \mathrm{Fin}\,n$ be a
set of coordinates, and let $k \ge 1$ be an integer satisfying $3k \le p\,\abs{U}$. Draw
a Bernoulli($p$) random restriction $\rho$ on $n$ variables, and let $J$ denote its set
of free variables. Then
\[
  \Pr\bigl[\,k \le \abs{U \cap J}\,\bigr] \;\ge\; \frac{1}{4}.
\]
\end{theorem}
